% =============================================================================
% 40-bab4.tex — BAB IV IMPLEMENTASI DAN TESTING
% =============================================================================
\chapter{IMPLEMENTASI DAN TESTING}
\label{chap:implementasi}

Bab ini memaparkan implementasi sistem berdasarkan perancangan yang telah diuraikan pada Bab~\ref{chap:analisa-perancangan}. Pembahasan mencakup lingkungan pengembangan, implementasi basis data, \textit{scoring engine}, API tRPC, sistem autentikasi, panel administrasi, generasi PDF dan pengiriman surel, pengujian unit, \textit{pipeline} CI/CD, serta \textit{deployment} ke lingkungan produksi.

% ─────────────────────────────────────────────────────────────────────────────
\section{Lingkungan Pengembangan}
\label{sec:lingkungan-dev}

Pengembangan Platform CHP dilakukan pada lingkungan lokal dengan spesifikasi perangkat lunak yang tercantum pada Tabel~\ref{tab:kebutuhan-software}. Manajer paket yang digunakan adalah pnpm versi 9, yang dipilih karena efisiensi penyimpanannya melalui mekanisme \textit{content-addressable store} yang menghindari duplikasi dependensi antar proyek.

Konfigurasi TypeScript menggunakan mode \texttt{strict: true} dengan opsi tambahan \texttt{noUncheckedIndexedAccess}, \texttt{noUnusedLocals}, dan \texttt{noUnusedParameters} untuk memaksimalkan deteksi kesalahan pada waktu kompilasi. Tiga \textit{path alias} dikonfigurasi untuk menyederhanakan impor modul: \texttt{@/*} untuk direktori \texttt{src/}, \texttt{@/server/*} untuk \texttt{src/server/}, dan \texttt{@/components/*} untuk \texttt{src/components/}.

Seluruh \textit{environment variable} yang diperlukan didokumentasikan dalam \textit{file} \texttt{.env.example} yang berisi 16 variabel konfigurasi, mencakup koneksi basis data (\texttt{DATABASE\_URL}), kunci autentikasi (\texttt{AUTH\_SECRET}, \texttt{ADMIN\_JWT\_SECRET}), konfigurasi layanan pihak ketiga (Upstash Redis, Sentry, Resend), dan kunci enkripsi (\texttt{ENCRYPTION\_KEY} sepanjang 64 karakter heksadesimal untuk AES-256-GCM).

% ─────────────────────────────────────────────────────────────────────────────
\section{Implementasi Basis Data}
\label{sec:impl-database}

Skema basis data yang dirancang pada Subbab~\ref{sec:perancangan-db} diimplementasikan menggunakan Drizzle ORM versi 1.0.0-beta.20 dengan \textit{driver} Neon HTTP (\texttt{@neondatabase/serverless}). Definisi skema diorganisasi ke dalam 14 \textit{file} TypeScript di dalam direktori \texttt{src/server/schema/}, dengan setiap \textit{file} bertanggung jawab atas satu domain.

Seluruh 20 tabel dan 7 \textit{enum} PostgreSQL berhasil di-\textit{deploy} ke lingkungan produksi Neon PostgreSQL. Migrasi skema dikelola melalui Drizzle Kit (\texttt{drizzle-kit}) yang menghasilkan \textit{file} SQL migrasi secara otomatis berdasarkan perubahan pada definisi skema TypeScript. Pendekatan ini memastikan bahwa setiap perubahan skema terdokumentasi dan dapat direproduksi.

Relasi antar-tabel diimplementasikan menggunakan \texttt{defineRelations()} dari Drizzle ORM v2, yang memisahkan definisi relasi dari definisi kolom. Pola ini memungkinkan kueri relasional (\textit{relational queries}) yang menghasilkan SQL \texttt{JOIN} secara otomatis tanpa memerlukan penulisan kueri manual.

Data awal (\textit{seed data}) untuk empat instrumen psikometri---PSS-10, GPIUS-2, SRS, dan SRQ-29---disiapkan melalui skrip \textit{seeding} yang memasukkan definisi tes, butir soal, opsi jawaban, dan rentang interpretasi. Skrip ini memastikan lingkungan pengembangan baru dapat langsung memiliki data yang cukup untuk pengujian.

% ─────────────────────────────────────────────────────────────────────────────
\section{Implementasi \textit{Scoring Engine}}
\label{sec:impl-scoring}

\textit{Scoring engine} diimplementasikan sebagai \textit{pipeline} fungsi murni (\textit{pure functions}) dalam tiga \textit{file} di direktori \texttt{src/server/scoring/}. Desain ini menjamin determinisme: untuk masukan yang sama, keluaran selalu identik, tanpa efek samping (\textit{side effects}) atau ketergantungan pada status eksternal.

\subsection{Fungsi \texttt{reverseScore}}

Fungsi \texttt{reverseScore(rawValue, options)} mengimplementasikan pembalikan skor secara agnostik skala menggunakan rumus:

\begin{equation}
  \text{reversedScore} = (\text{scaleMax} - \text{rawValue}) + \text{scaleMin}
\end{equation}

\noindent di mana \texttt{scaleMax} dan \texttt{scaleMin} diturunkan secara otomatis dari larik opsi yang diberikan. Pendekatan ini menghilangkan kebutuhan untuk mengonfigurasi batas skala secara manual, sehingga fungsi dapat digunakan untuk skala Likert dengan rentang apapun.

\subsection{Fungsi \texttt{computeScore}}

Fungsi \texttt{computeScore(input)} adalah inti dari \textit{scoring engine}. Fungsi ini menerima masukan berupa jawaban peserta beserta metadata soal, dan menghasilkan objek hasil yang mencakup:

\begin{enumerate}[leftmargin=2cm]
  \item Skor total (\texttt{totalScore}): penjumlahan seluruh skor butir soal setelah pembalikan dan pembobotan.
  \item Skor per dimensi (\texttt{dimensionScores}): pengelompokan skor berdasarkan kolom \texttt{dimension} pada setiap butir soal.
  \item Skor mentah per butir soal (\texttt{rawScores}): skor sebelum pembalikan untuk keperluan audit.
  \item Skor terhitung per butir soal (\texttt{computedScores}): skor setelah pembalikan dan pembobotan.
  \item Skor maksimum yang mungkin (\texttt{maxPossibleScore}): dihitung berdasarkan bobot dan opsi tertinggi.
\end{enumerate}

Untuk instrumen GPIUS-2 yang menggunakan metode \textit{Deficient Self-Regulation} (DSR), \texttt{computeScore} mengimplementasikan \textit{rollup} orde kedua: setelah menghitung skor per subskala, skor total dihitung sebagai penjumlahan subskala DSR yang merupakan gabungan dari beberapa dimensi primer.

\subsection{Fungsi \texttt{lookupInterpretation}}

Fungsi \texttt{lookupInterpretation(testId, totalScore, dimension?)} mengkueri tabel \texttt{result\_interpretations} untuk memetakan skor numerik ke label interpretasi yang dapat dibaca manusia. Fungsi ini mengembalikan label, deskripsi, rekomendasi, dan tingkat keparahan (\textit{severity level}). Pada kasus di mana tidak ditemukan rentang yang sesuai (\textit{interpretation miss}), fungsi mencatat peringatan ke Sentry tanpa melempar kesalahan, sehingga hasil asesmen tetap tersedia meskipun tanpa interpretasi.

\subsection{Fungsi Validasi}

Dua fungsi validasi---\texttt{validateAnswerValues} dan \texttt{validateCompleteness}---memastikan integritas data sebelum proses komputasi skor dimulai. Kedua fungsi bersifat murni (tanpa akses basis data), sehingga dapat diuji secara terisolasi. \texttt{validateAnswerValues} memeriksa bahwa setiap nilai jawaban berada dalam rentang opsi yang valid, sedangkan \texttt{validateCompleteness} memastikan semua butir soal wajib telah dijawab.

\subsection{Instrumen yang Diimplementasikan}

Empat instrumen psikometri berhasil dikonfigurasi dan divalidasi pada platform:

\begin{enumerate}[leftmargin=2cm]
  \item \textbf{PSS-10} (\textit{Perceived Stress Scale}): 10 butir soal, skala Likert 0--4, unidimensional, dengan 4 butir soal terbalik.
  \item \textbf{GPIUS-2} (\textit{Generalized Problematic Internet Use Scale}): 15 butir soal, skala Likert 1--5, 5 subskala dengan \textit{rollup} DSR orde kedua.
  \item \textbf{SRS} (\textit{Stress Response Scale}): 11 butir soal, skala Likert 1--6, 3 subskala.
  \item \textbf{SRQ-29} (\textit{Self-Reporting Questionnaire}): 29 butir soal, jawaban biner (Ya/Tidak), 4 kluster diagnostik.
\end{enumerate}

% ─────────────────────────────────────────────────────────────────────────────
\section{Implementasi API tRPC}
\label{sec:impl-trpc}

Seluruh 69 prosedur tRPC diimplementasikan dalam 17 \textit{file} prosedur yang diorganisasi ke dalam 16 \textit{router namespace}. Setiap prosedur mendefinisikan skema masukan menggunakan Zod dan dihubungkan ke salah satu dari lima tingkat \textit{middleware} autentikasi.

\subsection{Komposisi \textit{Router}}

\textit{File} \texttt{src/server/trpc/router.ts} berfungsi sebagai \textit{barrel file} yang mengomposisikan seluruh \textit{router} menjadi satu \textit{app router}. \textit{Router} dikelompokkan berdasarkan domain: \textit{router} publik (\texttt{publicTests}, \texttt{sessions}, \texttt{results}, \texttt{health}), \textit{router} pengguna (\texttt{profile}), dan \textit{router} admin (\texttt{admin}, \texttt{adminDashboard}, \texttt{adminTests}, \texttt{adminQuestions}, \texttt{adminScales}, \texttt{adminResults}, \texttt{reportRequests}, \texttt{adminAccounts}, \texttt{adminUserAccounts}, \texttt{adminAudit}, \texttt{adminProfile}).

\subsection{\textit{Structural Lock}}

Fitur keamanan data yang krusial pada prosedur admin adalah mekanisme \textit{structural lock}. Ketika sebuah instrumen tes telah memiliki sesi asesmen yang tercatat (diperiksa melalui fungsi \texttt{getSessionCount}), prosedur \texttt{updateTest}, \texttt{updateQuestion}, \texttt{deleteQuestion}, \texttt{reorderQuestions}, dan \texttt{updateRange} memblokir perubahan pada kolom yang sensitif terhadap penilaian---seperti \texttt{slug}, \texttt{scoringMethod}, bobot soal (\texttt{weight}), dan nilai opsi (\texttt{value}). Mekanisme ini mencegah invalidasi retroaktif terhadap data riset yang telah terkumpul.

\subsection{Operasi Idempoten}

Prosedur \texttt{saveProgress} menggunakan pola \texttt{ON CONFLICT ... DO UPDATE} (PostgreSQL \texttt{UPSERT}) pada \textit{constraint} unik komposit \texttt{(session\_id, question\_id)}, sehingga penyimpanan progres bersifat idempoten---pemanggilan berulang dengan data yang sama menghasilkan efek yang identik. Pendekatan ini menangani dengan aman skenario koneksi terputus di mana klien mungkin mengirim ulang jawaban yang sama.

% ─────────────────────────────────────────────────────────────────────────────
\section{Implementasi Autentikasi}
\label{sec:impl-auth}

Sistem autentikasi diimplementasikan sebagai dua subsistem yang sepenuhnya terpisah, masing-masing dengan mekanisme \textit{token}, penyimpanan \textit{cookie}, dan siklus hidup sesi yang independen.

\subsection{Autentikasi Pengguna Publik (Auth.js v5)}

Auth.js (NextAuth v5, versi 5.0.0-beta.30) digunakan untuk autentikasi pengguna publik dengan konfigurasi berikut:

\begin{enumerate}[leftmargin=2cm]
  \item \textbf{Strategi sesi}: JWT (\texttt{session: \{ strategy: "jwt" \}}).
  \item \textbf{\textit{Provider}}: \texttt{CredentialsProvider} dengan validasi surel dan kata sandi.
  \item \textbf{\textit{Password hashing}}: bcryptjs dengan \textit{cost factor} 12.
  \item \textbf{Adaptor}: \texttt{CHPDrizzleAdapter} kustom yang dipersiapkan untuk integrasi OAuth di masa mendatang.
  \item \textbf{\textit{Callbacks}}: \textit{callback} JWT menanamkan \texttt{userId} dan \texttt{role} ke dalam \textit{token}; \textit{callback} \texttt{session} meneruskan informasi ini ke konteks sesi.
\end{enumerate}

\subsection{Autentikasi Administrator (jose JWT)}

Subsistem autentikasi admin diimplementasikan secara mandiri menggunakan pustaka jose tanpa ketergantungan pada Auth.js:

\begin{enumerate}[leftmargin=2cm]
  \item \textbf{Algoritma}: HS256 dengan kunci rahasia \texttt{ADMIN\_JWT\_SECRET} (minimal 32 karakter).
  \item \textbf{\textit{Sliding window}}: setiap permintaan admin yang valid memperpanjang masa berlaku \textit{token} 30 menit dari waktu saat ini.
  \item \textbf{Batas keras}: \textit{token} tidak dapat diperpanjang melebihi 8 jam dari waktu penerbitan awal (\texttt{originalIat}), memaksa \textit{login} ulang untuk sesi yang sangat panjang.
  \item \textbf{Penyimpanan}: JWT disimpan dalam \textit{cookie} bernama \texttt{admin-token}, dikelola melalui modul \texttt{cookies.ts} yang menggunakan \texttt{ctx.resHeaders} (bukan \texttt{cookies()}) untuk kompatibilitas dengan konteks tRPC.
  \item \textbf{Revokasi}: prosedur \texttt{logout} menetapkan kolom \texttt{sessionInvalidatedAt} di basis data sebelum menghapus \textit{cookie}. \textit{Middleware} \texttt{adminProcedure} memeriksa kolom ini pada setiap permintaan, sehingga \textit{token} yang dicuri setelah \textit{logout} tidak dapat digunakan.
\end{enumerate}

\subsection{\textit{Rate Limiting} pada \textit{Login} Admin}

Modul \texttt{rate-limiter.ts} mengimplementasikan \textit{rate limiting in-memory} berlapis dua:

\begin{enumerate}[leftmargin=2cm]
  \item \textbf{Berbasis IP}: 5 kegagalan dalam jendela waktu 5 menit memicu penundaan (\textit{cooldown}) yang pulih secara otomatis setelah jendela berakhir.
  \item \textbf{Berbasis akun}: 10 kegagalan kumulatif mengunci akun secara permanen, memerlukan intervensi \texttt{super\_admin} melalui prosedur \texttt{unlockAdmin}. Mekanisme ini disimpan dalam \textit{in-memory} yang memadai untuk skala kurang dari 10 pengguna admin.
\end{enumerate}

% ─────────────────────────────────────────────────────────────────────────────
\section{Implementasi Panel Administrasi}
\label{sec:impl-admin}

Panel administrasi diakses melalui rute \texttt{/admin/(dashboard)/*} dan dilindungi oleh penjaga rute \texttt{proxy.ts} yang memeriksa keberadaan \textit{cookie} \texttt{admin-token}. Panel terdiri atas delapan modul utama:

\begin{enumerate}[leftmargin=2cm]
  \item \textbf{\textit{Dashboard}}: statistik agregat platform---jumlah asesmen selesai, sesi aktif, tingkat penyelesaian, tes terpopuler, dan hasil terbaru.
  \item \textbf{Manajemen Instrumen} (\texttt{adminTests}): CRUD tes dengan transisi status \texttt{draft} $\rightarrow$ \texttt{published} $\rightarrow$ \texttt{archived}, dilengkapi \textit{structural lock} untuk instrumen yang telah memiliki sesi.
  \item \textbf{Manajemen Soal} (\texttt{adminQuestions}): penambahan, penyuntingan, penghapusan, dan pengurutan ulang butir soal beserta opsi jawaban.
  \item \textbf{Konfigurasi Skala} (\texttt{adminScales}): pengelolaan rentang interpretasi per dimensi.
  \item \textbf{Data Hasil} (\texttt{adminResults}): tabel data dengan paginasi, filter, dan pengurutan; ekspor data dengan batas 5000 baris yang mencakup jawaban per butir soal.
  \item \textbf{Permintaan Laporan} (\texttt{reportRequests}): alur persetujuan bertahap (\textit{pending} $\rightarrow$ \textit{reviewed} $\rightarrow$ \textit{approved/rejected}) dengan dukungan persetujuan \textit{batch}.
  \item \textbf{Manajemen Akun Admin} (\texttt{adminAccounts}, khusus \texttt{super\_admin}): pembuatan akun, perubahan peran, aktivasi/deaktivasi, pembukaan kunci, dan invalidasi sesi paksa. Dilengkapi pengaman \textit{last super admin} yang mencegah penonaktifan \texttt{super\_admin} terakhir.
  \item \textbf{Jejak Audit} (\texttt{adminAudit}): log paginasi dengan filter cepat, pembatasan peran, dan resolusi label entitas. Setiap mutasi administratif menulis ke tabel \texttt{audit\_logs}.
\end{enumerate}

% ─────────────────────────────────────────────────────────────────────────────
\section{Implementasi Generasi PDF dan Pengiriman Surel}
\label{sec:impl-pdf-email}

Generasi laporan hasil asesmen menggunakan \texttt{@react-pdf/renderer} versi 4.5.1, yang memungkinkan pembuatan dokumen PDF menggunakan komponen React. Laporan yang dihasilkan mencakup informasi tes, skor total, skor per dimensi, label interpretasi, dan rekomendasi.

Pengiriman surel menggunakan layanan Resend versi 6.12.2, yang menyediakan API transaksional untuk pengiriman surel dari lingkungan \textit{serverless}. Alur lengkap---dari penerimaan permintaan melalui prosedur \texttt{approve} pada \textit{router} \texttt{reportRequests}, generasi PDF, hingga pengiriman surel---dieksekusi secara atomik dalam satu prosedur tRPC.

% ─────────────────────────────────────────────────────────────────────────────
\section{Pengujian Unit}
\label{sec:impl-unit-test}

Pengujian unit dilaksanakan menggunakan Vitest versi 4.1.5 dengan konfigurasi yang didefinisikan dalam \texttt{vitest.config.ts}. Total 235 kasus uji (\textit{test cases}) ditulis dalam 20 \textit{file} uji yang diorganisasi berdasarkan modul. Tabel~\ref{tab:inventaris-test} menyajikan inventaris lengkap.

\begin{longtable}{@{}p{6.5cm} c@{}}
  \caption{Inventaris \textit{File} Pengujian Unit}
  \label{tab:inventaris-test} \\
  \toprule
  \textbf{\textit{File} Uji} & \textbf{Jumlah Kasus Uji} \\
  \midrule
  \endfirsthead
  \toprule
  \textbf{\textit{File} Uji} & \textbf{Jumlah Kasus Uji} \\
  \midrule
  \endhead
  \bottomrule
  \endfoot
  \texttt{scoring/engine.test.ts} & $\sim$25 \\
  \texttt{scoring/reverseScore.test.ts} & $\sim$26 \\
  \texttt{scoring/validation.test.ts} & $\sim$21 \\
  \texttt{data/interpretations.test.ts} & $\sim$44 \\
  \texttt{data/srq29-questions.test.ts} & $\sim$8 \\
  \texttt{admin-questions/admin-questions.test.ts} & $\sim$18 \\
  \texttt{admin-accounts/admin-accounts.test.ts} & $\sim$16 \\
  \texttt{admin-scales/admin-scales.test.ts} & $\sim$12 \\
  \texttt{encryption.test.ts} & $\sim$12 \\
  \texttt{admin-auth/rate-limiter.test.ts} & $\sim$9 \\
  \texttt{reports/report-requests.test.ts} & $\sim$7 \\
  \texttt{admin-tests/admin-tests.test.ts} & $\sim$6 \\
  \texttt{admin-auth/jwt.test.ts} & $\sim$6 \\
  \texttt{admin-profile/admin-profile.test.ts} & $\sim$6 \\
  \texttt{admin-results/admin-results-delete.test.ts} & $\sim$5 \\
  \texttt{admin-audit/admin-audit.test.ts} & $\sim$4 \\
  \texttt{admin-accounts/admin-user-accounts.test.ts} & $\sim$4 \\
  \texttt{schema/admin-schema.test.ts} & $\sim$3 \\
  \texttt{schema/admin-tests-guards.test.ts} & $\sim$3 \\
  \texttt{smoke.test.ts} & $\sim$2 \\
  \midrule
  \textbf{Total} & \textbf{$\sim$235} \\
\end{longtable}

Modul \textit{scoring engine} diuji dengan 72 kasus uji yang tersebar dalam tiga \textit{file} (\texttt{engine.test.ts}, \texttt{reverseScore.test.ts}, \texttt{validation.test.ts}), mencakup skenario penjumlahan sumatif, pembalikan skor, pengelompokan dimensional, penanganan masukan kosong, validasi batas opsi, dan verifikasi kelengkapan jawaban.

% ─────────────────────────────────────────────────────────────────────────────
\section{Pengujian CI/CD}
\label{sec:impl-cicd}

\textit{Pipeline} CI/CD diimplementasikan dalam satu \textit{workflow} GitHub Actions (\texttt{.github/workflows/ci.yml}) yang diaktifkan pada setiap \textit{push} ke cabang \texttt{master} dan setiap \textit{pull request} yang menargetkan \texttt{master}. \textit{Pipeline} terdiri atas dua \textit{job} yang berjalan secara paralel:

\begin{enumerate}[leftmargin=2cm]
  \item \textbf{\textit{Job} \texttt{build-and-test}}: menggunakan Node.js 22 dan pnpm 9, menjalankan tiga \textit{gate} berurutan:
  \begin{enumerate}
    \item \texttt{pnpm exec eslint .} --- pemeriksaan \textit{linting}.
    \item \texttt{pnpm tsc --noEmit} --- pemeriksaan tipe.
    \item \texttt{pnpm build} --- \textit{build} produksi Next.js.
  \end{enumerate}
  \item \textbf{\textit{Job} \texttt{test}}: menggunakan Node.js 22 dan pnpm 9, menjalankan \texttt{pnpm test} untuk mengeksekusi seluruh \textit{test suite} Vitest.
\end{enumerate}

Kedua \textit{job} menggunakan \textit{caching} pnpm untuk mempercepat instalasi dependensi. Sebuah \textit{pull request} tidak dapat di-\textit{merge} apabila salah satu \textit{job} gagal.

% ─────────────────────────────────────────────────────────────────────────────
\section{\textit{Deployment}}
\label{sec:impl-deployment}

Platform CHP di-\textit{deploy} pada infrastruktur \textit{serverless} menggunakan dua layanan terkelola:

\begin{enumerate}[leftmargin=2cm]
  \item \textbf{Vercel} untuk komputasi: setiap \textit{push} ke cabang \texttt{master} secara otomatis memicu \textit{deployment} produksi. Setiap \textit{pull request} mendapatkan \textit{preview deployment} dengan URL unik untuk pengujian.
  \item \textbf{Neon PostgreSQL} untuk basis data: menggunakan \textit{pooled connection string} dengan fitur \textit{scale-to-zero} yang menghentikan komputasi basis data saat tidak ada permintaan, mengurangi biaya operasional.
\end{enumerate}

Seluruh 16 \textit{environment variable} dikonfigurasi melalui \textit{dashboard} Vercel dan tidak disimpan dalam repositori kode sumber. Kunci rahasia (\texttt{AUTH\_SECRET}, \texttt{ADMIN\_JWT\_SECRET}, \texttt{ENCRYPTION\_KEY}) dihasilkan menggunakan generator kriptografis yang aman dan disimpan secara terenkripsi oleh Vercel.

Pemantauan kesalahan (\textit{error monitoring}) menggunakan Sentry versi 10.50.0, yang terintegrasi baik pada sisi \textit{server} (\texttt{sentry.server.config.ts}) maupun sisi \textit{edge} (\texttt{sentry.edge.config.ts}). Setiap kesalahan yang tidak tertangani pada produksi secara otomatis dilaporkan ke \textit{dashboard} Sentry dengan konteks \textit{stack trace} dan metadata permintaan.

% ─────────────────────────────────────────────────────────────────────────────
\section{Pemenuhan Permintaan Dosen Pemilik Proyek}
\label{sec:pemenuhan-permintaan}

Selama pengembangan Platform Asesmen Digital CHP, dosen pemilik proyek menyampaikan 22 permintaan fungsional yang menjadi acuan pengembangan. Dari 22 permintaan tersebut, 18 permintaan telah tercapai dan 4 permintaan belum tercapai pada akhir periode kerja praktik. Tabel lengkap beserta status dan keterangan setiap butir permintaan disajikan pada Lampiran (lihat halaman~\ref{lamp:tabel-pemenuhan}).

\textbf{A.~Instrumen Asesmen.}
Permintaan penambahan dua instrumen baru---GPIUS-2 (\textit{Generalized Problematic Internet Use Scale 2}) dan SRS (\textit{Stress Response Scale})---berhasil dipenuhi dan ditambahkan ke katalog asesmen bersama instrumen yang telah ada sebelumnya, yaitu PSS-10 dan SRQ-29. Setiap instrumen dikonfigurasi secara lengkap meliputi butir soal, opsi jawaban, pembobotan, dimensi, dan rentang interpretasi, seluruhnya tervalidasi melalui 72 kasus uji pada \textit{scoring engine}.

\textbf{B.~Alur Asesmen Peserta.}
Alur asesmen peserta dilengkapi dengan visualisasi hasil secara \textit{real-time} yang ditampilkan segera setelah penyelesaian melalui komponen grafik Recharts (\textit{gauge}, \textit{radar}, diagram batang). Formulir demografis menyediakan pemilihan domisili bertingkat dua level (provinsi $\rightarrow$ kota/kabupaten) menggunakan dataset statis 38 provinsi yang bersumber dari klasifikasi BPS (Badan Pusat Statistik), dilengkapi opsi ``Lainnya (ketik manual)'' pada pemilihan provinsi maupun kota/kabupaten untuk mengakomodasi wilayah yang tidak terdaftar. Penyimpanan jawaban dilakukan secara otomatis pada setiap perubahan melalui prosedur \texttt{saveProgress} dengan pola \textit{upsert} idempoten, memungkinkan peserta melanjutkan sesi yang terputus.

\textbf{C.~Pelaporan dan Ekspor.}
Pengguna dapat mengajukan permintaan laporan melalui formulir tervalidasi; permintaan tersimpan pada tabel \texttt{report\_requests} untuk diproses administrator. Sistem menyediakan ekspor data dalam format CSV dan XLSX menggunakan pustaka SheetJS, dengan mekanisme \textit{filtered download} yang memastikan data yang diunduh bersifat dinamis mengikuti filter aktif (jenis kelamin, provinsi, kota, rentang usia, rentang skor, kategori, dan rentang tanggal). Setiap pengguna memiliki laporan komprehensif mencakup visualisasi data pribadi, analisis per-dimensi, dan interpretasi mendetail. Hasil mendalam dapat diunduh dalam format PDF (melalui \texttt{@react-pdf/renderer}), XLSX, dan CSV.

\textbf{D.~Dasbor dan Visualisasi Admin.}
Panel administrasi menyajikan dasbor dengan visualisasi agregat yang mendukung filter berdasarkan jenis instrumen dan rentang tanggal pada dasbor utama, serta filter lanjutan (pencarian nama, jenis kelamin, provinsi, kota, usia, skor, kategori, dan tanggal) pada halaman data hasil per instrumen. Dasbor antrean permintaan laporan memungkinkan administrator melihat, menyetujui, atau menolak permintaan; persetujuan memicu generasi PDF dan pengiriman surel secara atomik. Komponen grafik dilengkapi fitur unduh gambar menggunakan pustaka \texttt{html-to-image} (fungsi \texttt{toPng}) yang mengkonversi elemen DOM ke format PNG dengan resolusi 2$\times$ \textit{pixel ratio}. Kolom tabel hasil mendukung pengurutan dinamis \textit{ascending} dan \textit{descending} pada enam parameter (nama, jenis kelamin, usia, skor, kategori, dan tanggal tes) baik pada sisi server melalui kueri Drizzle ORM maupun pada sisi klien.

\textbf{E.~Manajemen dan Keamanan.}
Administrator dapat melakukan operasi CRUD penuh atas tes, butir soal, opsi jawaban, dan rentang interpretasi melalui modul manajemen asesmen yang dilengkapi mekanisme \textit{structural lock}. Setiap modifikasi administratif direkam dalam tabel \texttt{audit\_logs} mencakup aktor, aksi, entitas, serta nilai sebelum dan sesudah. Sistem menerapkan RBAC empat peran (\texttt{super\_admin}, \texttt{admin}, \texttt{psychiatrist}, \texttt{researcher}) dengan lima tingkat \textit{middleware} otorisasi: \texttt{publicProcedure} tanpa autentikasi, \texttt{protectedProcedure} untuk sesi Auth.js, \texttt{adminProcedure} untuk seluruh peran admin, \texttt{adminMutationProcedure} yang membatasi operasi tulis kepada \texttt{super\_admin} dan \texttt{admin}, serta \texttt{superAdminProcedure} yang eksklusif untuk \texttt{super\_admin}. Peran \texttt{psychiatrist} dan \texttt{researcher} memperoleh akses baca saja (\textit{read-only}) terhadap data hasil dan statistik. Modul \textit{Account Management} memungkinkan \textit{Super Admin} untuk membuat, memperbarui, dan mengatur hak akses akun administrator.

\subsection{Permintaan yang Belum Tercapai}
\label{subsec:belum-tercapai}

Empat permintaan belum tercapai pada akhir periode kerja praktik. Pertama, peringatan otomatis pada pertanyaan yang belum terisi (butir 2) belum diimplementasikan karena kompleksitas sinkronisasi status pengisian antara penyimpanan sisi-klien (\texttt{localStorage}) dan persistensi sisi-server secara simultan; sebagai mitigasi, tombol ``Kirim Asesmen'' dinonaktifkan secara \textit{default} dan hanya aktif setelah seluruh pertanyaan terjawab, sehingga pengiriman yang tidak lengkap tetap tercegah. Kedua, fitur \textit{import} data massal dalam format XLS/XLSX (butir 14) belum tercapai karena kompleksitas validasi skema terhadap data historis yang beragam format dan struktur, sehingga ditangguhkan ke iterasi berikutnya demi menjaga integritas data penelitian. Ketiga dan keempat, fitur \textit{Contact Us} untuk mengirim pesan ke surel administrator (butir 19) dan fitur \textit{Direct Mail Responder} untuk membalas pesan dari dasbor (butir 20) belum tercapai karena memerlukan integrasi layanan surel dua arah yang melampaui cakupan waktu kerja praktik.

Keempat permintaan yang belum tercapai telah diidentifikasi sebagai pekerjaan lanjutan dan didokumentasikan pada Subbab~\ref{sec:saran} sebagai saran pengembangan di masa mendatang.

Tabel pemenuhan selengkapnya yang mencakup seluruh 22 permintaan beserta status dan keterangan disajikan pada Lampiran (halaman~\ref{lamp:tabel-pemenuhan}).